\section{Related Works}\label{appendix:related_work}

\textbf{Web Agents:} Both proprietary and open-source communities have made significant strides in web agent development. Proprietary systems like DeepResearch~\citep{openai_deep_research} excel in complex web tasks but are hindered by closed architectures and inaccessible training data, limiting reproducibility and collaborative research. Open-source efforts, on the other hand, mainly focus on: data synthesis (e.g., data fuzzing in WebSailor and ASearcher), RL infrastructure, and algorithmic optimization (e.g., the specialized ARPO~\citep{dong2025arpo}). 
These advancements have propelled systems from addressing basic multi-hop question answering tasks to tackling more complex information-seeking challenges, such as the BrowseComp benchmark. Notable releases include WebSailor~\citep{li2025websailor}, WebShaper~\citep{tao2025webshaper}, ASearcher-QwQ-32B~\citep{gao2025asearcher}, WebExplorer-8B~\citep{liu2025webexplorer}, and DeepDive-32B~\citep{lu2025deepdiveadvancingdeepsearch}. Despite these achievements, open-source agents remain fundamentally limited by the constrained exploration capabilities of the ReAct paradigm~\citep{yao2023react}, highlighting the need for new paradigms.

\textbf{Context Management for Agents:} The most widely used approach for context management in LLM-based agents is ReAct's append-all-history strategy. While simple, this method leads to unbounded growth and rapid exhaustion, especially for complex queries. To address these issues, some methods introduce external components such as retrieval modules, e.g., A-MEM~\citep{xu2025mem}, MemOS~\citep{li2025memos_long}, and ReasoningBank~\citep{ouyang2025reasoningbank}, to structure context more effectively. However, these solutions add significant computational overhead, increase system complexity, and integrate loosely with the agent. More recent approaches, such as MEM1~\citep{zhou2025mem1}, MemAgent~\citep{yu2025memagent}, and MemSearcher~\citep{yuan2025memsearcher}, allow agents to manage context internally through RL. These methods innovate through \textbf{architectural modification}, introducing learnable memory tokens that require end-to-end training of a new agent from scratch, which limits their applicability to pre-existing agents.
In contrast, ReSum offers a lightweight \textbf{paradigmatic enhancement} to ReAct. This fundamental distinction yields key practical advantages including (1) \textbf{Training-free utility:} ReSum provides performance gains without any training, while methods like MEM1 can see performance drops in this setting (Table~\ref{tab:training_free_performance}); (2) \textbf{Data efficiency:} When training is required, ReSum-GRPO achieves significant improvements using only 1K samples (Table~\ref{tab:rl_performance}), reducing data cost; and (3) \textbf{Forward compatibility:} The decoupled summary tool can be independently improved, enhancing all ReSum-compatible agents without exhaustive retraining. 


\textbf{Distinction from World-model-augmented Agents:} While ReSum's summarization shares the high-level goal of providing structured guidance with world models proposed in \cite{chae2025web}, they address fundamentally different problems. World models focus on improving decision quality through dense, forward-looking planning at each step, whereas ReSum addresses context window exhaustion to enable decision continuation through sparse, backward-looking context compression. Therefore, the summary tool in ReSum functions as an \textbf{occasional context manager} rather than a \textbf{closely integrated planner}.
